\section{A scalar Gamma--Appell normalization of digamma residues}
\label{sec:scalar-gamma-normalizer}

The differential polynomials generated by a normalized Gamma quotient
have exactly the strict harmonic residues required for a single-depth
Laurent coefficient.  There is also a canonical scalar polynomial in
the digamma function that has the same residues up to a decaying
correction.  The distinction is useful: scalar powers are the original
finite-part integrands, whereas differential polynomials remove all
unwanted pole orders at once.

For $n\geq0$, let
\begin{equation}
 e_j(n)=\sum_{1\leq i_1<\cdots<i_j\leq n}
             \frac{1}{i_1\cdots i_j},
 \qquad e_0(n)=1,
 \qquad e_j(n)=0\quad(j>n).
 \label{eq:scalar-elementary-definition}
\end{equation}
Thus $e_j(n)$ is a strict finite harmonic sum, and
\begin{equation}
 E_n(t):=\sum_{j\geq0}e_j(n)t^j
   =\prod_{k=1}^{n}\left(1+\frac{t}{k}\right)
   =\frac{\Gamma(n+1+t)}{\Gamma(n+1)\Gamma(1+t)}.
 \label{eq:scalar-elementary-generator}
\end{equation}
Define the monic polynomials
\begin{equation}
 P_m(Y)=m![t^m]\Gamma(1-t)e^{tY},\qquad m\geq0.
 \label{eq:scalar-appell-definition}
\end{equation}
They satisfy $P_0=1$ and $P_m'(Y)=mP_{m-1}(Y)$, so they form an
Appell sequence.  For example,
\begin{equation}
 \begin{aligned}
 P_1(Y)&=Y+\gamma,\\
 P_2(Y)&=(Y+\gamma)^2+\zeta(2),\\
 P_3(Y)&=(Y+\gamma)^3+3\zeta(2)(Y+\gamma)+2\zeta(3).
 \end{aligned}
 \label{eq:scalar-first-polynomials}
\end{equation}
Every coefficient assertion below is at a fixed finite order; no
uniformity as $m\to\infty$ is required.
The generating function also gives an explicit finite recurrence:
\begin{equation}\label{eq:scalar-recurrence}
 P_{m+1}(Y)=(Y+\gamma)P_m(Y)
  +\sum_{r=1}^m\binom mr r!\zeta(r+1)P_{m-r}(Y).
\end{equation}
It follows by differentiating with respect to $t$ and using
$-\psi(1-t)=\gamma+\sum_{r\ge1}\zeta(r+1)t^r$.

\begin{lemma}[The universal cotangent residue]
\label{lem:scalar-cotangent-residue}
For every $t\in\mathbb C$,
\begin{equation}
 \operatorname*{Res}_{v=0}e^{-\pi t\cot(\pi v)}
       =-\frac{\sin(\pi t)}{\pi}.
 \label{eq:scalar-cotangent-residue}
\end{equation}
Here the residue is the coefficient of $v^{-1}$; the isolated
singularity may be essential.
\end{lemma}

\begin{proof}
Put $w=\pi\cot(\pi v)$.  Near $v=0$, this is an invertible change
between the local coordinate $v$ and the local coordinate $1/w$ at
infinity, since $w=v^{-1}+O(v)$.  Moreover,
\begin{equation}
 \frac{dw}{dv}=-(w^2+\pi^2),\qquad
 e^{-\pi t\cot(\pi v)}\,dv
       =-\frac{e^{-tw}}{w^2+\pi^2}\,dw.
 \label{eq:scalar-residue-change-variable}
\end{equation}
Residues of differentials are invariant under this local coordinate
change.  The two finite residues of the differential on the right are
$-e^{-i\pi t}/(2i\pi)$ at $w=i\pi$ and
$e^{i\pi t}/(2i\pi)$ at $w=-i\pi$.  Their sum is
$\sin(\pi t)/\pi$.  Applying the residue theorem on a circle
containing these two points, or equivalently taking the residue at
infinity, gives its negative, as asserted.  This argument applies to
every complex $t$; a pole rather than an essential singularity at
infinity is not needed for that contour calculation.
\end{proof}

\begin{theorem}[Scalar residue normalization and uniqueness]
\label{thm:scalar-gamma-normalization}
For $m\geq1$ and $n\geq0$, set
\begin{equation}
 r_{m,n}=\operatorname*{Res}_{x=-n}P_m(\psi(x)),
 \qquad
 \delta_{m,n}=r_{m,n}+m!e_{m-1}(n).
 \label{eq:scalar-defect-definition}
\end{equation}
Then $\delta_{1,n}=0$, and for each fixed $m\geq2$,
\begin{equation}
 \delta_{m,n}
   =O_m\!\left(\frac{(1+\log n)^{m-2}}{n}\right)
           \qquad(n\to\infty).
 \label{eq:scalar-defect-decay}
\end{equation}
In particular, the difference in
\eqref{eq:scalar-defect-definition} tends to zero.

Conversely, a monic polynomial $Q\in\mathbb C[Y]$ of degree $m$
satisfies
\begin{equation}
 \operatorname*{Res}_{x=-n}Q(\psi(x))+m!e_{m-1}(n)
           \longrightarrow0
 \label{eq:scalar-normalization-condition}
\end{equation}
if and only if $Q-P_m$ is constant.  Consequently
$P_m(Y)-P_m(0)$ is the unique such monic polynomial with zero
constant term.
\end{theorem}

\begin{proof}
Fix a circle $|v|=\rho$ with $0<\rho<1/2$.  The reflection formula
and the uniform large-argument expansion of the digamma function give
\begin{equation}
 \psi(-n+v)=\psi(n+1-v)-\pi\cot(\pi v)
           =L_n+K(v)+\eta_n(v),
 \label{eq:scalar-local-reflection}
\end{equation}
where $L_n=\log n$, $K(v)=-\pi\cot(\pi v)$, and
$\eta_n$ is holomorphic on $|v|\leq\rho$ with
$\sup|\eta_n|=O(n^{-1})$.  Only $n\geq1$ is needed for this
asymptotic argument.

Consider, coefficientwise near $t=0$,
\begin{equation}
 R_n(t)=\operatorname*{Res}_{v=0}e^{tK(v)+t\eta_n(v)}.
 \label{eq:scalar-residue-germ}
\end{equation}
Integration on the fixed circle makes $R_n(t)$ an entire function
of $t$.  By Lemma~\ref{lem:scalar-cotangent-residue},
\begin{equation}
 R_n(t)=-\frac{\sin(\pi t)}{\pi}+\Delta_n(t),
 \qquad
 \Delta_n(t)=\operatorname*{Res}_{v=0}
       e^{tK(v)}\bigl(e^{t\eta_n(v)}-1\bigr).
 \label{eq:scalar-residue-germ-error}
\end{equation}
Each fixed Taylor coefficient of $\Delta_n$ is $O(n^{-1})$:
this follows directly by bounding its contour integral on
$|v|=\rho$ for $t$ on a fixed small circle.  More precisely,
$\Delta_n(0)=0$ and
$\Delta_n'(0)=\operatorname*{Res}_{v=0}\eta_n(v)=0$.
Thus $\Delta_n$ starts with $t^2$, a fact that gives the exponent
$m-2$ in \eqref{eq:scalar-defect-decay}.

Taking coefficients in
\eqref{eq:scalar-appell-definition} and using
\eqref{eq:scalar-local-reflection}--\eqref{eq:scalar-residue-germ-error}
therefore yields, for $m\geq2$,
\begin{equation}
 \begin{aligned}
 r_{m,n}
 &=-m![t^m]\Gamma(1-t)e^{tL_n}
                         \frac{\sin(\pi t)}{\pi}
       +O_m\!\left(\frac{(1+L_n)^{m-2}}{n}\right)\\
 &=-m![t^{m-1}]\frac{e^{tL_n}}{\Gamma(1+t)}
       +O_m\!\left(\frac{(1+L_n)^{m-2}}{n}\right).
 \end{aligned}
 \label{eq:scalar-nondecaying-polynomial}
\end{equation}
The last equality is Euler's reflection identity
$\Gamma(1-t)\sin(\pi t)/\pi=t/\Gamma(1+t)$.

For comparison, uniformly for $t$ in a sufficiently small fixed disk,
\begin{equation}
 \frac{\Gamma(n+1+t)}{\Gamma(n+1)}
    =e^{tL_n}e^{u_n(t)},\qquad
 u_n(0)=0,\qquad [t^j]u_n(t)=O_j(n^{-1})\quad(j\geq1).
 \label{eq:scalar-gamma-ratio-error}
\end{equation}
For example, this follows by integrating
$\psi(n+1+s)-\log n=O(n^{-1})$ from $s=0$ to $s=t$.
Consequently the correction $e^{u_n(t)}-1$ starts with $t$ and
has each fixed coefficient $O(n^{-1})$.  By
\eqref{eq:scalar-elementary-generator}, for $d\geq1$,
\begin{equation}
 e_d(n)=[t^d]\frac{e^{tL_n}}{\Gamma(1+t)}
       +O_d\!\left(\frac{(1+L_n)^{d-1}}{n}\right).
 \label{eq:scalar-elementary-asymptotics}
\end{equation}
Equations \eqref{eq:scalar-nondecaying-polynomial} and
\eqref{eq:scalar-elementary-asymptotics}, with $d=m-1$, prove
\eqref{eq:scalar-defect-decay}.  For $m=1$ the assertion is
immediate from $P_1(\psi)=\psi+\gamma$, whose residue at every
nonpositive integer is $-1$.

It remains to prove necessity and uniqueness.  Let $Q$ satisfy
\eqref{eq:scalar-normalization-condition}.  If $Q-P_m$ is
nonconstant, write its degree as $j\geq1$ and its leading
coefficient as $a\ne0$.  For $j=1$, its residue is exactly
$-a$.  For $j\geq2$, expand the polynomial at
$L_n+K(v)+\eta_n(v)$ in \eqref{eq:scalar-local-reflection}.
The coefficient of $L_n^j$ is constant in $v$, hence contributes
no residue.  The coefficient of $L_n^{j-1}$ has residue
$ja\operatorname*{Res}_{v=0}(K+\eta_n)=-ja$.  All remaining
terms have degree at most $j-2$ in $L_n$ and bounded residue
coefficients.  Hence
\begin{equation}
 \operatorname*{Res}_{x=-n}(Q-P_m)(\psi(x))
      =-ja(\log n)^{j-1}
          +O\bigl((1+\log n)^{j-2}\bigr).
 \label{eq:scalar-uniqueness-leading-term}
\end{equation}
Neither this expression nor the nonzero constant for $j=1$ can
tend to zero, contradicting the two normalization conditions.
Thus $Q-P_m$ is constant.  Conversely, an additive constant has
zero residue and preserves the condition.
\end{proof}

The theorem establishes a uniqueness of functions, not an assertion
about algebraic independence of Euler's constant or zeta values.
It also explains why an unqualified uniqueness statement would be
false: no residue condition can determine an additive constant.

\begin{proposition}[Uniqueness of the complete Appell family]
\label{prop:scalar-appell-family-uniqueness}
Suppose $A(t)\in\mathbb C[[t]]$, $A(0)=1$, and
$Q_m(Y)=m![t^m]A(t)e^{tY}$.  If
\eqref{eq:scalar-normalization-condition} holds for every $m\geq1$,
then $A(t)=\Gamma(1-t)$ as formal power series.
\end{proposition}

\begin{proof}
By Theorem~\ref{thm:scalar-gamma-normalization},
$Q_m-P_m$ is constant for every $m$.  If
$A(t)-\Gamma(1-t)$ were nonzero, let $a_k t^k$, $a_k\ne0$,
be its first nonzero term; necessarily $k\geq1$.  The coefficient
of $t^{k+1}$ in
$(A(t)-\Gamma(1-t))e^{tY}$ contains the nonconstant term
$a_kY$.  This contradicts the conclusion for $m=k+1$.
\end{proof}

\subsection{An explicit finite formula for the defect}

Set $H_0^{(r)}=0$, write $H_n=H_n^{(1)}$, and define
\begin{equation}
 b_n=H_n-\gamma,\qquad
 a_{r,n}=H_n^{(r+1)}+(-1)^{r+1}\zeta(r+1)\quad(r\geq1).
 \label{eq:scalar-local-coefficients}
\end{equation}
The exact convergent local expansion is
\begin{equation}
 \psi(-n+v)=-\frac1v+b_n+\sum_{r\geq1}a_{r,n}v^r,
                         \qquad0<|v|<1.
 \label{eq:scalar-digamma-laurent}
\end{equation}
It follows by applying the recurrence to $\psi(v)$ and then using
the Taylor expansion of $\psi(1+v)$.

\begin{proposition}[Finite harmonic-polynomial defect formula]
\label{prop:scalar-finite-defect}
For every $m\geq2$ and $n\geq0$,
\begin{equation}
 \begin{aligned}
 r_{m,n}=m!\sum_{h=1}^{m}\frac{(-1)^h}{h!}
 [t^{m-h}v^{h-1}]\,
 \Gamma(1-t)
 \exp\left(tb_n+t\sum_{r=1}^{m-2}a_{r,n}v^r\right),\\
 \delta_{m,n}=r_{m,n}+m!e_{m-1}(n).
 \end{aligned}
 \label{eq:scalar-finite-defect-formula}
\end{equation}
The sum over $r$ is empty when $m=2$.  Formula
\eqref{eq:scalar-finite-defect-formula} requires only finitely
many coefficients of the displayed series; it is an exact
polynomial computation, valid also at $n=0$.
\end{proposition}

\begin{proof}
Substitute \eqref{eq:scalar-digamma-laurent} in
\eqref{eq:scalar-appell-definition} and expand
$e^{-t/v}=\sum_{h\geq0}(-t)^h/(h!v^h)$.  The residue from
the term indexed by $h$ is the coefficient of $v^{h-1}$ in
the holomorphic factor.  Only $1\leq h\leq m$ can contribute
to $[t^m]$.  Moreover, an occurrence of $a_{r,n}v^r$ consumes
at least one additional power of $t$, so whenever such a coefficient
occurs, $h\leq m-1$ and $r\leq h-1\leq m-2$.
All higher $a_{r,n}$ can therefore
be discarded before coefficient extraction.  This proves the
first equation; the second is the definition of the defect.
\end{proof}

For a transparent low-degree form, put
\begin{equation}
 A_j(n)=\zeta(j)-H_n^{(j)},\qquad j\geq2.
 \label{eq:scalar-harmonic-tails}
\end{equation}
The first defects are
\begin{equation}
 \begin{aligned}
 \delta_{1,n}&=\delta_{2,n}=\delta_{3,n}=0,\\
 \delta_{4,n}&=-4A_3(n),\\
 \delta_{5,n}&=-20H_nA_3(n)+5A_2(n)^2
                 -20\zeta(2)A_2(n)+25A_4(n).
 \end{aligned}
 \label{eq:scalar-low-degree-defects}
\end{equation}
These formulas can be checked without any infinite summation.
For example, writing $h=H_n$ and $h_j=H_n^{(j)}$, the finite
formula gives
\begin{equation}
 \begin{aligned}
 r_{4,n}&=-4h^3+12hh_2-4h_3-4\zeta(3),\\
 r_{5,n}&=-5h^4+30h^2h_2-20hh_3-20\zeta(3)h
           -10h_2^2+10\zeta(2)h_2\\
        &\hspace{1.5em}+5\zeta(2)^2+5h_4-25\zeta(4).
 \end{aligned}
 \label{eq:scalar-fourth-fifth-residues}
\end{equation}
Subtract the identities
$-24e_3(n)=-4h^3+12hh_2-8h_3$ and
$-120e_4(n)=-5h^4+30h^2h_2-15h_2^2-40hh_3+30h_4$.
The fifth defect then takes the form in
\eqref{eq:scalar-low-degree-defects} after
$\zeta(2)^2=5\zeta(4)/2$ is used.  In particular, scalar
normalization is already exact at the level of residues through
degree three.  At degrees four and five the remaining terms are
explicit decaying harmonic tails, rather than additional
nondecaying lower-depth data.

\subsection{Comparison with exact differential normalization}

In the notation of the Gamma-polygamma construction, let
\begin{equation}
 \mathcal C_m(x)=m![t^m]\frac{\Gamma(1-t)\Gamma(x)}{\Gamma(x-t)}.
 \label{eq:scalar-differential-comparison-definition}
\end{equation}
Taylor expansion of $\log\Gamma(x-t)$ gives the exact
coefficient identity
\begin{equation}
 \sum_{m\geq0}\frac{\mathcal C_m(x)}{m!}t^m
 =\Gamma(1-t)e^{t\psi(x)}
  \exp\left(\sum_{r\geq2}
       \frac{(-1)^{r+1}\psi^{(r-1)}(x)}{r!}t^r\right).
 \label{eq:scalar-differential-comparison-generator}
\end{equation}
Thus the higher derivatives provide an exact correction to the
scalar Appell family.  The first examples are
\begin{equation}
 \begin{aligned}
 \mathcal C_2&=P_2(\psi)-\psi',\\
 \mathcal C_3&=P_3(\psi)-3P_1(\psi)\psi'+\psi'',\\
 \mathcal C_4&=P_4(\psi)-6P_2(\psi)\psi'+4P_1(\psi)\psi''
                  -\psi'''+3(\psi')^2,\\
 \mathcal C_5&=P_5(\psi)-10P_3(\psi)\psi'
                   +10P_2(\psi)\psi''-5P_1(\psi)\psi'''\\
 &\qquad+15P_1(\psi)(\psi')^2-10\psi'\psi''+\psi''''.
 \end{aligned}
 \label{eq:scalar-differential-comparison-examples}
\end{equation}
Indeed the quotient itself has only simple possible poles in $x$,
and its exact residue is
\begin{equation}
 \operatorname*{Res}_{x=-n}
    \frac{\Gamma(1-t)\Gamma(x)}{\Gamma(x-t)}
       =-t\prod_{j=1}^{n}\left(1+\frac{t}{j}\right).
 \label{eq:scalar-differential-exact-residue}
\end{equation}
Coefficient extraction therefore yields
\begin{equation}
 \operatorname*{Res}_{x=-n} \mathcal C_m(x)=-m!e_{m-1}(n),
 \qquad
 \operatorname*{Res}_{x=-n}\bigl(P_m(\psi(x))-\mathcal C_m(x)\bigr)
          =\delta_{m,n}.
 \label{eq:scalar-exact-defect-comparison}
\end{equation}
The scalar theorem cancels all nondecaying residue data.  The
differential family also cancels the explicit defects and every
higher pole order.  These are precise functional statements;
they do not assert that the associated Laurent constants are
arithmetically independent or form a minimal period basis.

\subsection{A complete finite-part formula for inverse-Gamma derivatives}

There is an all-degree form of the cancellation which explains the
fifth-power identity before its raw-power endpoint corrections are
computed. Define
\begin{equation}\label{eq:Adefinition}
 \mathcal A_m(x)=(-1)^m\Gamma(x)\frac{d^m}{dx^m}\frac1{\Gamma(x)},
 \qquad
 \frac{\Gamma(x)}{\Gamma(x-t)}
       =\sum_{m\ge0}\mathcal A_m(x)\frac{t^m}{m!}.
\end{equation}
Then $\mathcal A_1=\psi$, $\mathcal A_2=\psi^2-\psi'$, and
$\mathcal A_{m+1}=\psi\mathcal A_m-\mathcal A_m'$.

\begin{corollary}[All-degree differential finite parts]
\label{cor:Aintegrals}
For $m\ge1$,
\begin{equation}\label{eq:Aintegrals}
 \boxed{\displaystyle
 \FP\int_0^1\mathcal A_m(x)\,dx
 =m![t^m]\frac{\sin(\pi t)}{\pi t}
   -m!\sum_{d=1}^{m-1}(-1)^{m-d-1}c_{m-d-1}\lambda_d.}
\end{equation}
The sum is empty for $m=1$. For odd $m$ the first term is zero;
for $m=2k$ it is $(-1)^k\pi^{2k}/(2k+1)$.
\end{corollary}
\begin{proof}
Multiplying the series for $\mathcal C_m$ by $1/\Gamma(1-t)$ gives
\[
 \mathcal A_m=\sum_{k=0}^m\frac{m!}{k!}
                         (-1)^{m-k}c_{m-k}\mathcal C_k.
\]
Finite parts are linear. The degree-zero integral is one, and
$\FP\int_0^1\mathcal C_1=\gamma$ because
$\int_\varepsilon^1\psi=-\log\Gamma(\varepsilon)$ has constant
term zero. For $k\ge2$ use Theorem~\ref{thm:Cintegral}.
The terms containing only $c_k$ combine to
\[
 m![t^m]\frac1{\Gamma(1-t)\Gamma(1+t)}
     =m![t^m]\frac{\sin(\pi t)}{\pi t},
\]
and the remaining terms are exactly the displayed finite sum.
\end{proof}

In degree five this reads
\begin{equation}\label{eq:A5height}
 \begin{aligned}
 \FP\int_0^1\mathcal A_5(x)\,dx
   ={}&-120\xi+120\gamma\theta
      -60(\gamma^2-\zeta(2))\eta\\
      &+20(\gamma^3-3\gamma\zeta(2)+2\zeta(3))\lambda_1.
 \end{aligned}
\end{equation}
The next section proves exactly how $\mathcal A_5$ differs, after
integration with the specified finite parts, from the scalar power
$\psi^5$. The resulting additional mixed integral is an analytic
quantity, not an omitted integration constant.
